\section*{Capítulo \ref{ch:g8:combustion-and-fuels} --- A combustão: combustíveis, produtos e gases de efeito estufa}

\begin{solution}{exo:g8:combustion-and-fuels:1}
\ce{C + O2 -> CO2}.
\end{solution}

\begin{solution}{exo:g8:combustion-and-fuels:2}
Com oxigênio suficiente, todo o carbono vira dióxido de carbono
(completa). Com pouco oxigênio, uma parte vira monóxido de carbono ou
fuligem (incompleta).
\end{solution}

\begin{solution}{exo:g8:combustion-and-fuels:3}
Ele não tem cor, nem cheiro, nem gosto: ninguém o percebe, e mesmo assim
ele é tóxico.
\end{solution}

\begin{solution}{exo:g8:combustion-and-fuels:4}
Fósseis: carvão mineral, gás natural, gasolina. Renováveis: lenha,
biogás, etanol de beterraba.
\end{solution}

\begin{solution}{exo:g8:combustion-and-fuels:5}
\ce{C3H8 + 5O2 -> 3CO2 + 4H2O}.
\end{solution}

\begin{solution}{exo:g8:combustion-and-fuels:6}
Carbono: 2 \ce{CO2}; hidrogênio: 3 \ce{H2O}; oxigênio à direita $4 + 3 =
7$, dos quais 1 vem do etanol, então 3 \ce{O2}:
\ce{C2H6O + 3O2 -> 2CO2 + 3H2O}.
\end{solution}

\begin{solution}{exo:g8:combustion-and-fuels:7}
Fuligem, isto é, carbono. A \omterm{def:g4:what-a-fire-needs:burning}{combustão} é incompleta: falta oxigênio à
chama.
\end{solution}

\begin{solution}{exo:g8:combustion-and-fuels:8}
Cerca de 339 ppm em 1980 e cerca de 370 ppm em 2000: um aumento de cerca de 30 ppm.
\end{solution}

\begin{solution}{exo:g8:combustion-and-fuels:9}
O carbono da lenha foi tirado do \omterm{def:g4:air-a-mixture-of-gases:air}{ar} pela árvore há poucos anos ou
décadas: queimá-la devolve esse dióxido de carbono. O carbono do carvão
mineral estava guardado no subsolo havia milhões de anos: queimá-lo
acrescenta dióxido de carbono novo ao \omterm{def:g4:air-a-mixture-of-gases:air}{ar}.
\end{solution}

\begin{solution}{exo:g8:combustion-and-fuels:10}
\ce{2C + O2 -> 2CO}.
\end{solution}

\begin{solution}{exo:g8:combustion-and-fuels:11}
$(427 - 316)/316 = 111/316 \approx 0.35$: cerca de \qty{35}{\percent}.
\end{solution}

\begin{solution}{exo:g8:combustion-and-fuels:12}
No começo há oxigênio suficiente: \ce{CH4 + 2O2 -> CO2 + 2H2O}. À medida
que o oxigênio do quarto fechado é consumido, a chama passa a queimar
com falta de \omterm{def:g4:air-a-mixture-of-gases:air}{ar}: \ce{2CH4 + 3O2 -> 2CO + 4H2O}. O monóxido de carbono,
sem cheiro e tóxico, vai se acumulando no quarto.
\end{solution}

\begin{solution}{pb:g8:combustion-and-fuels:1}
\textbf{1.} \ce{C3H8 + 5O2 -> 3CO2 + 4H2O}.

\textbf{2.} À esquerda: 3 C, 8 H, 10 O. À direita: 3 C, $4 \times 2 = 8$ H,
$3 \times 2 + 4 = 10$ O. Balanceada.

\textbf{3.} \ce{2C3H8 + 7O2 -> 6CO + 8H2O}.

\textbf{4.} A \omterm{def:g8:combustion-and-fuels:complete}{combustão completa}: 5 \omterm{def:g7:atoms-and-molecules:molecule}{moléculas} de oxigênio por \omterm{def:g7:atoms-and-molecules:molecule}{molécula}
de propano, contra $7/2 = 3.5$ na incompleta.

\textbf{5.} O \omterm{def:g4:air-a-mixture-of-gases:air}{ar} da barraca fechada logo fica com falta de oxigênio; a
chama não recebe mais o suficiente e \omterm{def:g4:what-a-fire-needs:burning}{queima} de forma incompleta,
amarela.

\textbf{6.} O monóxido de carbono: sem cor, sem cheiro (e sem gosto).

\textbf{7.} GHS02, a chama: extremamente inflamável. GHS04, o cilindro
de gás: um gás guardado sob pressão.

\textbf{8.} Por exemplo: usá-lo só ao \omterm{def:g4:air-a-mixture-of-gases:air}{ar} livre, nunca numa barraca nem
num cômodo fechado; mantê-lo longe de tudo o que pode queimar; trocar o
cartucho longe de qualquer chama.

\textbf{9.} $132 \div 44 = 3$ vezes mais pesado.

\textbf{10.} Do oxigênio do \omterm{def:g4:air-a-mixture-of-gases:air}{ar}, que acaba no dióxido de carbono (e na
água).

\textbf{11.} Um \omterm{def:g8:combustion-and-fuels:fossil}{combustível fóssil} (ele vem do gás natural ou do
petróleo): queimá-lo acrescenta dióxido de carbono ao \omterm{def:g4:air-a-mixture-of-gases:air}{ar}.

\textbf{12.} $450 \times 3 = \qty{1350}{g} = \qty{1.35}{kg}$ de dióxido
de carbono.
\end{solution}
